\begin{lemma}[The model field of (1.3) of the thesis inhabits (Q1) of the thesis--(Q2) of the thesis at $\varepsilon=0$ on $1<p<2$]
Let $n\ge1$ be an integer and let $1<p<2$.  Then the field
\[
A(x,z,h)=|h|^{p-2}h ,\qquad A(x,z,0)=0 ,
\]
of (1.3) of the thesis --- written $(|h|^2)^{\frac{p-2}2}h$, so that the value at
$h=0$ is the one the convention of \noderef{EllipticStructureQ12} assigns ---
satisfies
\[
\mathrm{EllipticStructureQ12}\bigl(n,p,\ \varepsilon=0,\ \gamma_0=p-1,\
\gamma_1=3-p,\ \Omega=\R^n,\ A\bigr).
\]

\noindent\textbf{Why this node exists.}  The Note for the formalization requires
at l.~258 that a cited result ``must not be vacuous (each must be a true
statement for the operators to which it is applied)''.  The same demand applies
a fortiori to the structure bundle every such citation is conditioned on: if
$\mathrm{EllipticStructureQ12}$ were satisfied by no field, Theorems
Theorem~1 of the thesis--Theorem~2 of the thesis and everything downstream would be vacuously true and
the development would certify nothing.  This node discharges that demand for the
one field the thesis names, (1.3) of the thesis, at the one parameter value where the
risk is real: $\varepsilon=0$ with $1<p<2$, the singular range the thesis is
about, where the bundle as printed in the thesis is \emph{contradictory} and
\noderef{EllipticStructureQ12} therefore restricts clauses (4)--(6) to
$D(\varepsilon)$.  That restriction is justified in that node's commentary; this
node is the machine-checked half of the justification, and because it is
Lean-closed it is re-checked whenever the bundle is edited --- which is what
caught the two Carath\'eodory clauses (1) and (2) when they were added.

The parameters are the ones \noderef{EllipticStructureQ12} records:
$\gamma_0=p-1$ and $\gamma_1=3-p$, so $\gamma_1/\gamma_0=(3-p)/(p-1)$ depends
only on $p$, as the constants of Theorems~Theorem~1 of the thesis and~Theorem~2 of the thesis require.
$\gamma_1=3-p$ rather than $1$ because the tablet's (Q2) is the componentwise
bound $|D_{h_j}A^i|\le\gamma_1(\varepsilon+|h|^2)^{\frac p2-1}$, not a bound on
the quadratic form.  The statement is given on $\Omega=\R^n$ because the field
is $x$-independent; a bundle on any $\Omega\subseteq\R^n$ follows, since every
clause of \noderef{EllipticStructureQ12} is a condition quantified over
$x\in\Omega$ or over all $x$.
\end{lemma}
\begin{proof}
Write $W=|h|^2$ and $A(h)=W^{\frac{p-2}2}h$ for the field, which does not
depend on $x$ or on $z$.  Note $p-1>0$ and $3-p>0$, so the two parameters are
admissible, and $D(0)=\R^n\setminus\{0\}=\{h:W\ne0\}$.

\emph{Clauses (1) and (2), the Carath\'eodory pair.}  For fixed $z$ and $h$ the
map $x\mapsto A(h)$ is constant, hence Borel measurable.  For clause (2), fix
$x$; the map $(z,h)\mapsto A(h)$ is the composition of the projection
$(z,h)\mapsto h$ with $A$, so it suffices that $A$ is continuous on $\R^n$,
which is clause (3).

\emph{Clause (3), continuity in $h$ on all of $\R^n$.}  Away from the origin
$h\mapsto W$ is smooth and positive, so $W^{\frac{p-2}2}$ is smooth there by
smoothness of $t\mapsto t^{\frac{p-2}2}$ on $(0,\infty)$, and $A$ is smooth as a
product of a scalar and the identity.  At the origin, $|A(h)|=|h|^{p-1}$: for
$h\ne0$,
\[
|A(h)|=\bigl|W^{\frac{p-2}2}\bigr|\,|h|=|h|^{p-2}|h|=|h|^{p-1},
\]
using $W^{\frac{p-2}2}=(|h|^2)^{\frac{p-2}2}=|h|^{p-2}$ for $|h|>0$ and
$W^{\frac{p-2}2}>0$; and at $h=0$ both sides are $0$, because $p-1\ne0$ makes
$0^{p-1}=0$ under the real-power convention.  Since $p-1>0$, the right-hand side
tends to $0$ as $h\to0$, so $A(h)\to0=A(0)$ by the squeeze.  Hence $A$ is
continuous on $\R^n$.

\emph{Clause (4), $C^1$ on $D(0)$.}  At $h$ with $W\ne0$ the map $y\mapsto|y|^2$
is $C^\infty$ and does not vanish at $h$, so $y\mapsto W(y)^{\frac{p-2}2}$ is
$C^1$ at $h$, and $A$ is $C^1$ at $h$ as the product of that scalar with the
identity.  Its derivative is
\begin{equation}
D_hA(h)\xi=W^{\frac{p-2}2}\xi+(p-2)\,W^{\frac{p-2}2-1}(h,\xi)\,h ,
\tag{$\ast$}
\end{equation}
by the product and chain rules, using $D(|y|^2)\xi=2(h,\xi)$ at $y=h$ and
$\frac{p-2}2\cdot2=p-2$.

\emph{Clause (5), (Q1).}  Fix $h\in D(0)$, so $W>0$, and fix $\xi$.  Pairing
$(\ast)$ with $\xi$ and using $(\,\cdot\,,\cdot\,)$-bilinearity and
$(\xi,\xi)=|\xi|^2$,
\[
\sum_{i,j}D_{h_j}A^i(h)\xi^i\xi^j=(D_hA(h)\xi,\xi)
=W^{\frac{p-2}2}|\xi|^2+(p-2)\,W^{\frac{p-2}2-1}(h,\xi)^2 .
\]
By Cauchy--Schwarz $(h,\xi)^2\le W|\xi|^2$, and $2-p>0$, so
\[
(2-p)W^{\frac{p-2}2-1}\bigl(W|\xi|^2-(h,\xi)^2\bigr)\ge0 .
\]
Expanding and using $W^{\frac{p-2}2-1}\cdot W=W^{\frac{p-2}2}$ (valid since
$W>0$) this says exactly
\[
W^{\frac{p-2}2}|\xi|^2+(p-2)W^{\frac{p-2}2-1}(h,\xi)^2
\;\ge\;(p-1)\,W^{\frac{p-2}2}|\xi|^2 ,
\]
which is (Q1) with $\gamma_0=p-1$, since $\varepsilon=0$ gives
$(\varepsilon+|h|^2)^{\frac p2-1}=W^{\frac{p-2}2}$.  The bound holds for
\emph{every} $x$, hence a.e.

\emph{Clause (6), (Q2).}  Fix $h\in D(0)$ and $i,j$.  Taking $\xi=e_j$ in
$(\ast)$ and reading off the $i$-th component, $(h,e_j)=h_j$ gives
\[
D_{h_j}A^i(h)=W^{\frac{p-2}2}\delta_{ij}+(p-2)\,W^{\frac{p-2}2-1}h_jh_i .
\]
Here $|\delta_{ij}|\le1$ and $|h_jh_i|\le|h|^2=W$, the latter from
$|h_k|\le|h|$ for each $k$ --- which is Cauchy--Schwarz against the unit vector
$e_k$.  Writing $d=\delta_{ij}$, $u=h_jh_i$ and using
$W^{\frac{p-2}2-1}W=W^{\frac{p-2}2}>0$ together with $2-p>0$,
\[
\bigl|W^{\frac{p-2}2}d+(p-2)W^{\frac{p-2}2-1}u\bigr|
\le W^{\frac{p-2}2}+(2-p)W^{\frac{p-2}2}=(3-p)W^{\frac{p-2}2},
\]
which is (Q2) with $\gamma_1=3-p$.  Again this holds for every $x$, hence a.e.
\end{proof}
